% Ad Familiares 13.38 --- headnote
% ad-familiares-13-38, 46 BC, Rome

\textit{Cicero to Manius Acilius Glabrio, proconsul of
Sicily, written from Rome in 46 BC (the manuscript
dateline: \textit{Scr. Romae, ut videtur, a. 708 (46)}).
The seventh surviving letter of the Acilius cluster
(Fam.\ 13.30--39). The beneficiary is Lucius Bruttius, a
young Roman knight with property in Sicily, whose
father had been a friend of Cicero's from his
quaestorship at Lilybaeum (75 BC). Bruttius himself is
at Rome; what needs the proconsul's attention is the
estate left under the management of agents in the
province.}

\textit{Notable for the explicit appeal to Cicero's
Sicilian quaestorship as the root of the connection ---
the office, more than two decades behind him, that had
given him a permanent client-base on the island and
furnished the material for the prosecution of Verres in
70 BC. The letter belongs to a recurring pattern in the
\textit{Ad Familiares}: an Italian principal at home, an
estate and a roster of named or unnamed agents in the
province, and a proconsul asked to extend protection to
the second through regard for the first. The closing
formula \textit{id quod ei recepi} (``what I undertook
to him for'') is a delicate touch: Cicero has already
promised Bruttius that the letter will produce its
effect, and is putting the proconsul on notice of the
promise.}
